% ECONOMICS_PROBLEM: {"id": "O33-23", "title": "Selection and attraction of knowledge-growth traveling waves", "jel_code": "O33", "source_id": "OP-0212"}
% FIRST_POSTED: 2026-10-09
% ATTEMPT_STATUS: unresolved
% ENTRY_KIND: research_attempt
\documentclass[11pt]{article}
\usepackage[T1]{fontenc}
\usepackage[utf8]{inputenc}
\usepackage{amsmath,amssymb,amsthm}
\usepackage[margin=1in]{geometry}
\usepackage{hyperref}
\newtheorem{theorem}{Theorem}
\newtheorem{proposition}[theorem]{Proposition}
\newtheorem{lemma}[theorem]{Lemma}
\newtheorem{corollary}[theorem]{Corollary}
\newcommand{\R}{\mathbb R}
\newcommand{\E}{\mathbb E}
\newcommand{\T}{\mathbb T}
\newcommand{\Pp}{\mathbb P}
\newcommand{\dd}{\mathrm d}
\title{Knowledge-growth wave selection:\\finite values, state sensitivity and vanishing upper-tail search}
\author{GPT 6 Astra Ultra --- Version 2}
\date{9 October 2026}
\begin{document}
\maketitle
\noindent\textbf{Problem O33-23. Status: unresolved.} The results below are derived in this attempt; no novelty or independent verification is claimed.

\section{Target and scope}
The target is attraction, after median centering, to a Lucas--Moll balanced-growth wave, including selection of its speed. We retain the infinite-horizon value boundary condition. Wave existence does not answer this question. The known wave theorem is Porretta--Rossi, Theorem 2.2 \cite{PR}; the time-dependent question is discussed in \cite[Section 3.3]{CD}. We retain the necessary speed and integrability restrictions from Version 1. Version 2 adds a finite-value theorem, a uniform productivity-sensitivity bound, and a quantitative upper-tail search restriction. These are restrictions on admissible solutions, not an attraction theorem.

Write $a=\max_{s\in[0,1]}\alpha(s)$ and $q(t,x)=\alpha(s_*(t,x))$. The density equation is
\begin{equation}\label{density}
 f_t=\nu f_{xx}+f(x)\int_{-\infty}^x q(y)f(y)\dd y-q(x)f(x)\int_x^\infty f(y)\dd y,
 \qquad 0\le q\le a.
\end{equation}
Time arguments inside integrals are suppressed. We assume a positive classical probability-density solution. Whenever integration by parts with an unbounded weight is used, assume the corresponding exponential moment is locally bounded in time and that the weak formulation extends to that weight by cutoff; sufficient conditions are integrability of $e^{\lambda x}(f+|f_x|)$ and vanishing weighted boundary terms. This additional moment regularity is stated, not inferred just from pointwise classical regularity.

\section{Tail and moment equations}
\begin{proposition}\label{moment}
Let $w(t,x)=\int_x^\infty f(t,y)\dd y$ and $A(t,x)=\int_{-\infty}^x q(t,y)f(t,y)\dd y$. If the flux at infinity vanishes, then
\[
 w_t=\nu w_{xx}+Aw,\qquad 0\le A\le a(1-w).
\]
For $\lambda>0$ with the stated moment regularity, $M_\lambda(t)=\int e^{\lambda x}f(t,x)\dd x$ satisfies
\begin{align}\label{momentidentity}
 M_\lambda'(t)&=\nu\lambda^2 M_\lambda(t)+I_\lambda(t),\\
 I_\lambda(t)&=\int_{x<y}q(t,x)f(t,x)f(t,y)(e^{\lambda y}-e^{\lambda x})\dd x\dd y,
 \qquad 0\le I_\lambda\le aM_\lambda.\nonumber
\end{align}
In particular, if $b(t)$ is a median,
\[
 b(t)\le\frac{\log(2M_\lambda(0))}{\lambda}
       +\left(\nu\lambda+\frac a\lambda\right)t.
\]
Every existing limit $b(t)/t\to c$ therefore satisfies
$c\le\inf_{\lambda\in\Lambda}(\nu\lambda+a/\lambda)$,
where $\Lambda$ is any set of positive exponents for which these moment hypotheses hold.
\end{proposition}
\begin{proof}
Integrate \eqref{density} over $(x,\infty)$. In the reaction terms all ordered meetings with both states above $x$ cancel. Only meetings starting below $x$ and ending above $x$ remain, giving $Aw$. The diffusion contribution is $-\nu f_x=\nu w_{xx}$. The bound on $A$ follows from $q\le a$ and $\int_{-\infty}^x f=1-w$.

Multiplying \eqref{density} by $e^{\lambda x}$, integrating the diffusion twice by parts, and interchanging the nonnegative ordered-meeting integrals gives \eqref{momentidentity}. Its upper bound follows from
\[
 I_\lambda\le a\int e^{\lambda y}f(t,y)\int_{-\infty}^y f(t,x)\dd x\dd y\le aM_\lambda.
\]
Gronwall's inequality yields $M_\lambda(t)\le M_\lambda(0)e^{(\nu\lambda^2+a)t}$. Since at least half the mass lies above the median, $M_\lambda(t)\ge e^{\lambda b(t)}/2$. These two inequalities prove the assertion, including its infimum version.
\end{proof}

The hypotheses of the question give a first exponential moment initially. They do not supply every exponent in $\Lambda$. In particular, replacing the infimum by $2\sqrt{\nu a}$ needs the exponent $\sqrt{a/\nu}$ to be available. Even when this upper bound holds, it does not identify the selected wave, because $A$ depends on the endogenous optimal policy.

\section{Restrictions on a possible limit wave}
\begin{theorem}\label{waves}
Let $(\varphi,S,c)$ be a positive classical traveling density satisfying the density equation, $\int\varphi=1$, and the flux and moment conditions above. Put
\[
 W(z)=\int_z^\infty\varphi(y)\dd y,\quad
 B(z)=\int_{-\infty}^z\alpha(S(y))\varphi(y)\dd y,\quad
 \beta=B(+\infty).
\]
Then
\begin{equation}\label{tailode}
 \nu W''+cW'+BW=0.
\end{equation}
If $\beta>0$ and the first exponential moment is finite, then $c>\nu$ and
$c\ge2\sqrt{\nu\beta}$. More generally, whenever $M_\lambda(\varphi)<\infty$ with valid weighted integrations,
\begin{equation}\label{speedmoment}
 c=\nu\lambda+\frac{I_\lambda(\varphi,S)}{\lambda M_\lambda(\varphi)}
 \le\nu\lambda+\frac a\lambda.
\end{equation}
If in addition $h(z)=-W'(z)/W(z)$ converges to a finite positive $\ell$ as $z\to\infty$, then
\begin{equation}\label{dispersion}
 c=\nu\ell+\frac\beta\ell.
\end{equation}
Finite first moment implies $\ell\ge1$; when $\ell>1$ it guarantees first-moment integrability at the right tail, whereas $\ell=1$ alone decides neither integrability nor nonintegrability.
\end{theorem}
\begin{proof}
The tail equation from Proposition~\ref{moment} with $w(t,x)=W(x-ct)$ is \eqref{tailode}. A traveling density has $M_\lambda(t)=e^{\lambda ct}M_\lambda(\varphi)$, giving \eqref{speedmoment}. When $\beta>0$, the set where $\alpha(S(x))\varphi(x)>0$ has positive measure. Positivity of $\varphi$ gives positive mass strictly above every such $x$, so $I_1>0$ and $c>\nu$.

Set $Y(z)=e^{cz/(2\nu)}W(z)>0$. Then
\[
 Y''+\left(\frac{B(z)}\nu-\frac{c^2}{4\nu^2}\right)Y=0.
\]
If $c^2<4\nu\beta$, the coefficient is at least some $\eta>0$ on a half-line. No positive function solves $Y''\le-\eta Y$ there: if $Y'$ becomes negative, concavity forces a zero in finite distance; if $Y'$ stays nonnegative, $Y$ stays bounded below by a positive constant and $Y''$ is uniformly negative, forcing $Y'$ to become negative. Hence $c^2\ge4\nu\beta$. The already proved $c>0$ gives the claimed lower bound.

Substitution of $h=-W'/W$ in \eqref{tailode} gives
$h'=h^2-(c/\nu)h+B/\nu$. If $h\to\ell$, its right-hand side has a limit. That limit must be zero, since a derivative with a nonzero limit is incompatible with convergence of its primitive. This proves \eqref{dispersion}. Finally,
\[
 \int_0^\infty e^z\varphi(z)\dd z=W(0)+\int_0^\infty e^zW(z)\dd z
\]
when finite, by integration by parts or Tonelli. If $\ell<1$, choose $\ell<r<1$; eventually $h\le r$, so $W(z)\ge C e^{-rz}$ and the last integral diverges. If $\ell>1$, choose $1<r<\ell$; eventually $h\ge r$, so $W(z)\le C e^{-rz}$ and it converges. At $\ell=1$, the tails $e^{-z}(1+z)^{-2}$ and $e^{-z}$ show the two alternatives. These last tails illustrate the integrability test; they are not asserted to be wave solutions.
\end{proof}

The published wave analysis gives stronger restrictions under its full value-profile hypotheses; Theorem~\ref{waves} is deliberately only a necessary test from the density and moment equations. It does not reprove the existence result of \cite{PR}.

\begin{proposition}[The square-root search rule]
Suppose $\alpha(s)=a\sqrt{s}$ and the wave value $V$ is nondecreasing with finite meeting integral. Define
\[
 Q(z)=e^{-z}\int_z^\infty[V(y)-V(z)]\varphi(y)\dd y.
\]
Then the maximizing search fraction is uniquely
\[
 S(z)=\min\{1,a^2Q(z)^2/4\}.
\]
If $V$ is the finite infinite-horizon value and $\rho>\nu$, then
$V(z)\ge e^z/(\rho-\nu)$.
\end{proposition}
\begin{proof}
Divide the objective by $e^z$ and set $r=\sqrt{s}\in[0,1]$. The maximized expression becomes $1-r^2+aQ(z)r$, a strictly concave quadratic whose maximizer is $\min\{1,aQ(z)/2\}$ because $Q\ge0$. Squaring proves the formula. Choosing zero search forever is admissible and leaves only diffusion; its expected discounted production is $e^z\int_0^\infty e^{-(\rho-\nu)t}\dd t$. Taking the supremum proves the lower bound.
\end{proof}


\section{Finite values and productivity sensitivity: progress beyond Version 1}
Version 1 assumed that the infinite-horizon value was finite and proved
only its no-search lower bound. We now verify finiteness under the
catalogue's discount condition and control its dependence on the initial
state. These arguments are conditional on a density flow satisfying the
moment assumptions of Proposition~\ref{moment}; they do not construct a
solution of the coupled density--value system.

Throughout this section assume $a>0$, as in the catalogue's concrete
square-root wave regime. If $a=0$, there are no meetings: the value is
exactly $e^x/(\rho-\nu)$ for $\rho>\nu$, and $\beta=0$; the strict
conclusion $\beta<a$ below is not asserted in that case.

Fix a starting time $t$. Write $A_s\in[0,1]$ for a control at elapsed time
$s$ to distinguish it from the constant rate bound $a$. Work on a filtered
space supporting a Brownian motion $B$, a rate-$a$ Poisson process,
and independent uniform marks $(U,V)$ at its points, independent of $B$.
Let $F_r^{-1}$ be a measurable quantile function of the given density
$f(r,\cdot)$. At a potential meeting at elapsed time $s$, accept the meeting
if $U\leq\alpha(A_s)/a$ and, if accepted, replace the log state by
$\max\{X_{s-},F_{t+s}^{-1}(V)\}$. Between meetings the state has diffusion
$\sqrt{2\nu}\,dB_s$. Controls are predictable for this driving filtration,
possibly with additional independent private randomization; in particular,
they do not anticipate the current marks. This is the usual thinning
construction of the controlled meeting intensity. There are finitely many
potential meetings on every bounded interval, so the construction is
pathwise well defined. The same construction covers progressively
measurable controls after the usual predictable representative at
Poisson times; it does not allow a control to inspect a meeting mark
before choosing whether to search.

For this admissible class define
\[
 J_{t,x}(A)=\E\int_0^\infty e^{-\rho s}(1-A_s)e^{X_s}\,ds,
 \qquad v(t,x)=\sup_A J_{t,x}(A).
\]
The controls are indexed by the common driving filtration, rather than
restricted to a prescribed feedback function. Thus a control chosen for
one starting state can be used for another starting state: any auxiliary
state trajectory it uses is itself a causal function of the same driving
noise. This point is needed when taking suprema in the coupling proof.
The argument includes arbitrary randomized admissible deviations in this
formulation.

\begin{lemma}[Productivity under bounded meeting rates]\label{productivity}
Suppose $M_1(t+s)\leq M_1(t)e^{(\nu+a)s}$ for $s\geq0$ and $M_1(t)<\infty$.
For every admissible control and $s\geq0$,
\[
 \E e^{X_s}\leq e^{\nu s}e^x
       +a\int_0^s e^{\nu(s-r)}M_1(t+r)\,dr.                 \tag{20}
\]
If two processes start at $x\leq x'$ and use the same control and all the
same driving noise, then $X_s\leq X'_s$ pathwise and
\[
 0\leq\E(e^{X'_s}-e^{X_s})\leq e^{\nu s}(e^{x'}-e^x).       \tag{21}
\]
\end{lemma}
\begin{proof}
Let $Z_s=e^{X_s}$ and $\sigma=\sqrt{2\nu}$. Between meetings,
$dZ_s=\nu Z_s\,ds+\sigma Z_s\,dB_s$.
The continuous multiplicative factor is $L_s=e^{\sigma B_s}$, with
$\E L_s=e^{\nu s}$. At a potential meeting at time $r$ the jump
$\Delta Z_r$ is nonnegative and at most the mark
$e^{F_{t+r}^{-1}(V)}$, even when the meeting is rejected.
Variation of constants on the finitely many meeting intervals gives
\[
 Z_s=e^xL_s+\sum_{r\leq s}\frac{L_s}{L_r}\Delta Z_r
 \leq e^xL_s+
       \sum_{r\leq s}\frac{L_s}{L_r}e^{F_{t+r}^{-1}(V_r)}.
\]
The sum is over potential meetings, so its right-hand side is independent
of acceptance decisions. Tonelli's theorem and the rate-$a$ marked
Poisson expectation formula apply to its nonnegative terms.
The Brownian increment $B_s-B_r$ is independent of the meeting time and
mark and has exponential expectation $e^{\nu(s-r)}$. Hence the expected
sum is $a\int_0^s e^{\nu(s-r)}M_1(t+r)dr$, proving (20) and its
finiteness. The Poisson expectation formula here can also be obtained by
conditioning on the number and times of its points and integrating the
independent uniform marks.

For the pair of starting states, order is preserved both by their common
Brownian translation in log space and by the map $z\mapsto\max\{z,y\}$.
In productivity coordinates the latter map is
$z\mapsto\max\{z,e^y\}$; for $z'\geq z$ it satisfies
\[
 0\leq\max\{z',e^y\}-\max\{z,e^y\}\leq z'-z.
\]
Acceptance is the same for the coupled processes because the control is
the same. Consequently $D_s=Z'_s-Z_s$ follows the common geometric
Brownian factor between meetings and has nonpositive jumps, while staying
nonnegative. Iterating across the finite set of meetings yields
$0\leq D_s\leq(e^{x'}-e^x)L_s$. Taking expectations proves (21).
\end{proof}

\begin{theorem}[A finite value and a uniform increment bound]\label{finitevalue}
Assume the moment regularity of Proposition~\ref{moment} holds on every
finite time interval, $0\leq\alpha\leq a$, and $\rho>\nu+a$.
Then the value defined above is finite and obeys
\begin{equation}\label{valuebounds}
 \frac{e^x}{\rho-\nu}\leq v(t,x)
 \leq\frac{e^x}{\rho-\nu}
       +\frac{aM_1(t)}{(\rho-\nu)(\rho-\nu-a)}.
\end{equation}
For every $h\geq0$,
\begin{equation}\label{increment}
 0\leq v(t,x+h)-v(t,x)
       \leq\frac{e^{x+h}-e^x}{\rho-\nu}.
\end{equation}
In particular $v(t,\cdot)$ is locally Lipschitz. Where it is differentiable,
$0\leq v_x(t,x)\leq e^x/(\rho-\nu)$; for the classical values in the
catalogue this holds at every point.
\end{theorem}
\begin{proof}
Proposition~\ref{moment}, applied from time $t$, gives the moment bound
required in Lemma~\ref{productivity}. Since $0\leq1-A_s\leq1$,
(20), Tonelli's theorem, and then the moment bound give
\begin{align*}
 J_{t,x}(A)
 &\leq\frac{e^x}{\rho-\nu}
 +a\int_0^\infty e^{-\rho s}\int_0^s
                    e^{\nu(s-r)}M_1(t+r)\,dr\,ds\\
 &=\frac{e^x}{\rho-\nu}
 +\frac a{\rho-\nu}\int_0^\infty e^{-\rho r}M_1(t+r)\,dr\\
 &\leq\frac{e^x}{\rho-\nu}
 +\frac{aM_1(t)}{(\rho-\nu)(\rho-\nu-a)}.
\end{align*}
The bound is uniform in the control. Zero search gives the lower bound
in \eqref{valuebounds}. Both denominators are positive by assumption.

For a fixed control use the ordered coupling in (21). The payoff
increment is nonnegative, and its expectation is at most
\[
 \int_0^\infty e^{-\rho s}\E D_s\,ds
 \leq\frac{e^{x+h}-e^x}{\rho-\nu}.
\]
Thus $J_{t,x}(A)\leq J_{t,x+h}(A)\leq J_{t,x}(A)+C_{x,h}$ for
every control in the same admissible class. Taking suprema proves
\eqref{increment}; no optimizer or feedback regularity is required.
Applying that bound on a compact interval proves local Lipschitz
continuity, and taking difference quotients proves the derivative bound.
\end{proof}

\begin{corollary}[Search vanishes in the upper tail]\label{vanishingsearch}
Under Theorem~\ref{finitevalue}, suppose $\alpha(s)=a\sqrt{s}$.
For the value-based maximizing search selector define
\[
 Q_t(x)=e^{-x}\int_x^\infty[v(t,y)-v(t,x)]f(t,y)\,dy.
\]
Then
\begin{align}
 0\leq Q_t(x)&\leq\frac1{\rho-\nu}
       \int_x^\infty(e^{y-x}-1)f(t,y)\,dy,\label{Qtail}\\
 s_*(t,x)&=\min\{1,a^2Q_t(x)^2/4\},\qquad
           \lim_{x\to\infty}s_*(t,x)=0.\nonumber
\end{align}
In particular every positive finite-first-moment wave in this regime
satisfies $S(z)\to0$ as $z\to\infty$ and
\[
 \beta=\int_\R a\sqrt{S(z)}\varphi(z)\,dz<a.
\]
The limit is for each fixed time or wave; no uniform-in-time centered-tail
estimate is asserted.
\end{corollary}
\begin{proof}
Apply \eqref{increment} to each $y\geq x$ and integrate. This proves
\eqref{Qtail}, including finiteness of its meeting integral. The right
side is at most
$e^{-x}(\rho-\nu)^{-1}\int_x^\infty e^yf(t,y)dy$, which tends to zero
as $x\to\infty$ by the finite first moment. Monotonicity of the value
makes $Q_t\geq0$. Maximizing in $\sqrt{s}$ as in the earlier search-rule
proposition gives the stated unique selector and its limit.
For a wave evaluate at time zero, where the value is $V$ and density is
$\varphi$. Choose $R$ sufficiently large that $S(z)\leq1/4$ for
$z\geq R$. Positivity of $\varphi$ gives
$\int_R^\infty\varphi>0$, and hence
\[
 \beta\leq a\int_{-\infty}^R\varphi
           +\frac a2\int_R^\infty\varphi<a.
\]
\end{proof}

\section{A compactness obstruction and the first gap}
\begin{proposition}
Convergence of median-centered densities in $L^1$ does not imply convergence of their first exponential moments, even for strictly positive smooth densities.
\end{proposition}
\begin{proof}
Let $\varphi$ be the standard Gaussian density and $\varphi_n(x)=(1-e^{-n})\varphi(x)+e^{-n}\varphi(x-2n)$. Then $\|\varphi_n-\varphi\|_1\le2e^{-n}$, but
$\int e^x\varphi_n(x)\dd x=e^{1/2}(1-e^{-n}+e^n)\to\infty$.
Let $b_n$ be the unique median of $\varphi_n$. Uniform convergence of the cumulative distribution functions and strict increase of the Gaussian cumulative function imply $b_n\to0$. Translation continuity in $L^1$ gives $\varphi_n(b_n+\cdot)\to\varphi$ in $L^1$, while their moments equal $e^{-b_n}\int e^x\varphi_n(x)\dd x\to\infty$.
\end{proof}
Thus neither the exact moment identity nor the nonlocal value integral can be passed to a proposed median-frame limit using the stated density topology alone. A candidate route must first prove suitable uniform integrability of the centered tails and control of the backward value field, then establish attraction and distinguish competing wave branches. No such estimate, attraction theorem, or basin characterization is proved here. The new finite-value and sensitivity theorems verify additional admissibility properties and rule out persistently maximal search in the upper tail. They do not supply the missing moving-frame estimates or dynamical attraction; the complete target remains unresolved.

\section{Completion Estimate}
20\%

This is a heuristic assessment of the full catalogue target, not a probability that the proofs are correct. The version proves value finiteness, state sensitivity and upper-tail search decay in addition to the retained speed and moment restrictions. It does not prove attraction to a wave, identify a basin of initial densities, or select a speed or profile; those are the principal unresolved parts of the target.

\begin{thebibliography}{9}
\bibitem{CD} P. Cardaliaguet and F. Delarue, \emph{Selected topics in mean field games}, Proc. ICM 2022, vol. 5, 3660--3703, EMS Press (2023), Section 3.3. \url{https://doi.org/10.4171/ICM2022/17}.
\bibitem{PR} A. Porretta and L. Rossi, \emph{Traveling waves for a nonlocal KPP equation and mean-field game models of knowledge diffusion}, arXiv:2010.10828v1 (2020), system (2), Theorem 2.2 and its following remarks. \url{https://arxiv.org/pdf/2010.10828v1}.
\end{thebibliography}

\end{document}
